\section{Related Work}
\label{sec:related-work}

\textbf{Automated harness optimization.} Prior work has extensively studied the automated optimization of prompts for LLM-based systems~\citep{zhou2023ape,pryzant2023protegi,yuksekgonul2025textgrad,khattab2024dspy,agrawal2025gepa}. Subsequent work extends automated optimization to agents, using execution trajectories and task feedback to improve their memory, skills, and workflows~\citep{zhao2024expel,zhang2025ace,ni2026trace2skill,hu2024adas,zhang2024aflow}. Meta-Harness and AHE further extend automated optimization to executable harness code, using trajectories and evaluation feedback to guide iterative revisions~\citep{lee2026metaharness,lin2026agentic}.

Building on this setting, offline methods revise harnesses using reference answers or task-completion feedback from training tasks, aiming to improve held-out performance~\citep{park2026autosaddler,zhang2026selfharness,lin2026agentic}. However, an evaluation finds unreliable gains without unit-test feedback and limited transfer from reusable harnesses, but reports gains from Harness Scaling, which alternates task attempts and harness revisions at test time~\citep{wang2026harnessevolution}.

Following this task-specific direction, JIT-Agent synthesizes a tailored harness for each task before execution begins~\citep{zhang2026jita}. TTHE further explores harness optimization without reference answers, using trajectories and execution-derived feedback to revise executable harnesses on unlabeled test batches~\citep{nie2026tthe}. Continual Harness moves harness optimization into ongoing task execution through an outer loop that periodically invokes a refiner while preserving the task state~\citep{karten2026continualharness}. In \boxbox{}, the task-solving agent decides when to allocate its remaining budget to full-harness revision, while a separate kernel manages runtime hot reloading and preserves execution context.

\textbf{Long-Horizon and Open-Ended Agent Evaluation.} Existing agent benchmarks cover diverse capabilities but commonly evaluate tasks with explicit end states, reporting task success, pass rates, or final artifact quality~\citep{mialon2023gaia, liu2024agentbench, jimenez2024swebench, xie2024osworld, merrill2026terminalbench, sun2026agentslastexam}. Recent iterative optimization benchmarks evaluate solution quality but still use reference implementations, human performance, or reference outputs for task construction or scoring~\citep{chan2025mlebench,lyu2026mlsbench,imajuku2025alebench,xu2026autolab,mang2025frontiercs,zhu2026edgebench}. Closer to our setting, Frontier-Eng and FrontierSWE support iterative optimization of engineering designs and software projects, with tasks curated from established problems and expert contributions~\citep{chi2026frontiereng,chu2026frontierswe}.

In contrast, \perfrace{} constructs long-running performance optimization tasks from real-world software projects and workloads without requiring reference optimizations. Its construction pipeline provides a scalable way to continually source new tasks from public software projects.

\textbf{Benchmarks for Performance Optimization.} \citet{shypula2024learning} introduce PIE for evaluating LLMs' ability to make performance-improving edits to single-file competitive programming solutions while preserving functional correctness. Subsequent benchmarks broaden the optimization domains to GPU kernels, general-purpose numerical programs, and algorithmic redesign, while retaining function- or file-level optimization tasks~\citep{du2024mercury,huang2024effibench,liu2024evalperf,waghjale2024ecco,ouyang2025kernelbench,press2025algotune}.

GSO and SWE-Perf extend evaluation to repository-level performance optimization, deriving tasks and human reference optimizations from GitHub commits and pull requests~\citep{shetty2025gso,he2025sweperf}. SWE-fficiency and FormulaCode further incorporate a targeted workload per task and multiple community-maintained workloads, respectively~\citep{ma2025swefficiency,sehgal2026formulacode}.

These repository-level benchmarks construct tasks around a targeted performance test, whereas real performance engineering starts from an end-to-end workload and discovers bottlenecks through profiling. \perfrace{} adopts this workload-first formulation, leaving agents to identify and implement correctness-preserving improvements without a known target patch.
